% =============================================================
% Retinal Vessel Segmentation (DRIVE) -- Literature Review, M1
% IEEE Conference Report -- 2 to 3 pages, no figures
%
% Compile: pdflatex literature_review.tex   (run twice)
% =============================================================

\documentclass[conference]{IEEEtran}
\IEEEoverridecommandlockouts

\usepackage{cite}
\usepackage{amsmath,amssymb,amsfonts}
\usepackage{graphicx}
\usepackage{textcomp}
\usepackage{xcolor}
\usepackage{booktabs}
\usepackage{array}
\usepackage{multirow}
\usepackage{float}
\usepackage[hyphens]{url}
\usepackage{hyperref}
\hypersetup{colorlinks=true,linkcolor=blue,citecolor=blue,urlcolor=blue}
\usepackage[font=small,labelfont=bf,skip=4pt]{caption}

\def\BibTeX{{\rm B\kern-.05em{\sc i\kern-.025em b}\kern-.08em
    T\kern-.1667em\lower.7ex\hbox{E}\kern-.125emX}}

\begin{document}

% =============================================================
\title{Deep Learning for Retinal Vessel Segmentation on DRIVE:\\
A Focused Literature Review and Base-Paper Selection}

\author{%
\setlength{\tabcolsep}{3pt}%
\begin{tabular}{@{}*{4}{>{\centering\arraybackslash}p{0.235\textwidth}}@{}}
% NOTE: do NOT use \IEEEauthorblockN inside this tabular -- the macro
% contains its own line break, which would end the row and stack the
% names into column 1. Plain \normalsize matches its typography exactly.
{\normalsize Tejmul Movin}                                             &
{\normalsize Eranki Sai Vikas}                                         &
{\normalsize Meenaksh Singhania}                                       &
{\normalsize Akula Jitendranath}                                       \\[3pt]
{\small\textit{Dept.\ of Computer Science}}                            &
{\small\textit{Dept.\ of Computer Science}}                            &
{\small\textit{Dept.\ of Computer Science}}                            &
{\small\textit{Dept.\ of Computer Science}}                            \\
{\small Newton School of Technology}                                   &
{\small Newton School of Technology}                                   &
{\small Newton School of Technology}                                   &
{\small Newton School of Technology}                                   \\
{\small Rishihood University}                                          &
{\small Rishihood University}                                          &
{\small Rishihood University}                                          &
{\small Rishihood University}                                          \\
{\small Sonipat, Haryana, India}                                       &
{\small Sonipat, Haryana, India}                                       &
{\small Sonipat, Haryana, India}                                       &
{\small Sonipat, Haryana, India}                                       \\
{\small sansar.t23csai@}                                               &
{\small eranki.v23csai@}                                               &
{\small meenaksh.s23csai@}                                             &
{\small akula.j23csai@}                                                \\
{\small nst.rishihood.edu.in}                                          &
{\small nst.rishihood.edu.in}                                          &
{\small nst.rishihood.edu.in}                                          &
{\small nst.rishihood.edu.in}
\end{tabular}
}

\maketitle

% =============================================================
\begin{abstract}
This review studies ten papers published from 2021 to 2026 on marking blood
vessel pixels in retinal fundus photographs. We compare their network designs,
training schemes, metrics, model sizes and, most importantly for us, whether
their code can actually be run. Reported $F_1$ scores sit between
\textbf{82.06\%} and \textbf{86.44\%}, but those numbers cannot be lined up
directly, because the papers mask, threshold and split their data differently.
Three gaps stand out: four of the ten share no code, progress has slowed to
about four points of $F_1$ in five years, and every method is trained pixel by
pixel even though a vessel tree is only useful when it stays joined up. We
select \textbf{SGL} (MICCAI 2021) as the base paper, because its repository
already holds the DRIVE images, the labels and the trained weights, and because
its loss can be swapped from the command line. After reproducing its baseline
we will test one controlled change: replacing the pixel-wise loss with the
centre-line clDice loss. No clinical claim is made in this work.
\end{abstract}

\begin{IEEEkeywords}
retinal vessel segmentation, DRIVE, U-Net, reproducibility, topology-aware
loss, medical image segmentation.
\end{IEEEkeywords}

% =============================================================
\section{Scope and Review Method}

The task is binary pixel classification. The input is one colour photograph of
the retina, and the output is a mask that marks every pixel as vessel or
background. Telling arteries apart from veins, finding lesions, grading disease
severity, and any work on OCT volumes all sit outside this project.

\subsection{Search and Selection}

We searched using terms such as \textit{retinal vessel segmentation},
\textit{DRIVE}, \textit{U-Net}, \textit{topology loss} and \textit{lightweight
vessel segmentation}. From the results we kept ten papers, published in 2021 or
later, that solve this exact task and report scores on the DRIVE test set.
Papers evaluated only on other vessel datasets were dropped. One of them,
Rolling-UNet (AAAI 2024), was removed at this stage because it tests on
CHASE\_DB1 rather than DRIVE. This is a focused course review, not a formal
systematic review.

\subsection{Benchmark and Metrics}

The benchmark is DRIVE, released by Staal \textit{et al.} in 2004. There are 40
colour images in it, each $565 \times 584$ pixels, gathered during a Dutch
screening drive for diabetic eye disease. The people who built the dataset also
fixed the split for everyone: 20 images to train on and 20 to test on. Every
image carries a field-of-view (FOV) mask, and every test image was marked by
two people, of whom the first is treated as ground truth. Because the split is
fixed, two groups reporting DRIVE scores are at least measuring the same
images.

Only about 8 to 10 per cent of pixels are vessel. A model that predicts pure
background therefore scores above 0.90 accuracy while finding nothing, so
accuracy is close to useless here. We compare $F_1$ (also called Dice),
sensitivity, specificity, AUC and IoU, and we treat $F_1$ as the main score.

% =============================================================
\section{U-Net Variants and Training-Scheme Methods}

Nearly all recent work starts from U-Net (Ronneberger \textit{et al.}, 2015)
and changes one part of it. The first group changes the network so that thin
vessels are not lost.

FR-UNet~\cite{frunet} argues that the repeated shrinking inside U-Net throws
away one-pixel vessels, and that enlarging the image later cannot bring them
back. It therefore keeps a full-size path from start to finish and passes
features between sizes. It then adds a two-threshold clean-up step: pixels with
a high score are kept, and pixels with a middling score are kept only when they
touch a pixel that was already kept. Its repository is one of only two in this
review that ships trained weights. MDFI-Net~\cite{mdfinet} worries about the
same loss, and answers it with deformable convolutions inside a pulse-coupled
block, so weak vessels are still visible deeper down. Its reported accuracy
runs from 97.91 to 98.16 per cent over three datasets. No code was released.

The second group changes the training instead of the network.
SGL~\cite{sgl} makes the point that the limit may be the labels. Many thin
vessels in DRIVE were never marked by the annotator, so a model trained on them
learns to miss the same vessels. SGL trains $k$ models on different parts of
the data, combines them into cleaner labels, and trains again on those. It also
supplies a tool that rubs out vessel labels on purpose, so robustness to bad
labels can be measured instead of assumed. MRC-Net~\cite{mrcnet} blends
features taken at several scales, lets information travel in both directions,
and trains against a discriminator. The point of that last part is to reward
the model for whole regions rather than for isolated pixels.

% =============================================================
\section{Topology, Attention, and Lightweight Models}

A third group looks at the shape of the vessel itself. DSCNet~\cite{dscnet}
lets the sampling points of a convolution curve along the vessel instead of
sitting on a rigid square grid, and pairs this with a loss that charges the
model directly for broken shapes. It reports 82.06 per cent Dice on DRIVE and
has the most active repository in this review. This is the closest work to our
own idea, but it changes the convolution and the loss at the same time, so the
effect of each part cannot be separated.

A fourth group widens the context the model can see.
SA-UNetv2~\cite{saunetv2} builds on the authors' own earlier SA-UNet. It adds
attention at each skip link, and it trains with a loss that pairs weighted
cross-entropy with a correlation term. The result is 82.82 per cent $F_1$ from
just 0.26 million weights, and about a second of work on a plain CPU.
MM-UNet~\cite{mmunet} trades its pointwise convolutions for Mamba-style
state-space layers, and gains 1.64 per cent $F_1$ on DRIVE.

The last group cares about cost. LVS-Net~\cite{lvsnet} and
LFRA-Net~\cite{lfranet} report Dice scores of 86.44 and 84.28 per cent using
0.71 and 0.17 million parameters. Neuron programming~\cite{neuronprog} is the
outlier of the group. It hunts for useful neuron types instead of a better
layout, then trains a small helper network to guess good settings so the search
stays affordable.

% =============================================================
\section{Main Findings and Research Gaps}

The reported range of 82.06 to 86.44 per cent $F_1$ is not a fair ranking. Some
papers score inside the FOV mask and some do not, thresholds are picked
differently, and patch sizes vary, so a gap of one point can come from the
protocol alone. A 2024 benchmarking study by Fadugba \textit{et al.} found that
these models fall off sharply once they meet a dataset they were never tuned
for, which hints that the published gaps are smaller than they look.

Three gaps stand out. First, four of the ten papers release no code, so nobody
can check their numbers. Second, scores have moved only about four points in
five years while model sizes have fallen roughly ten times, so recent progress
is mostly in efficiency rather than accuracy. Third, and most useful for us,
every method above is trained with a loss that scores one pixel at a time, yet
no paper reports a connectivity score. This matters because losing a single
pixel in the middle of a thin branch barely changes Dice, while it cuts the
branch in two and spoils any later measurement of branching or curviness.
DSCNet~\cite{dscnet} is the only reviewed work that attacks this directly, and
it does so together with a new convolution.

We will study this gap without changing the network at all. That keeps the
experiment small, clear, and possible on a free GPU.

% =============================================================
\section{Base Paper and Proposed Improvement}

\subsection{Why SGL}

We selected SGL~\cite{sgl} for reproduction; Table~\ref{tab:checklist} gives
our check. Six of the ten papers share code, and we compared those six on how
easily they run. SGL won for four reasons. Its repository already contains the
DRIVE images, the labels, the generated pseudo-labels and the trained weights,
so there is nothing to download and no folder to rearrange. A single command
loads the published weights and prints the test scores, which tells us within
minutes whether our setup is correct. It needs only ordinary PyTorch and NumPy,
with one small fix expected because newer PyTorch changed the default of
\texttt{torch.load}. And its loss is selected from the command line, with a
loss folder that already holds cross-entropy, Dice and focal loss, so our
change is one new file and one flag rather than surgery on the training loop.

Our first choice had been FR-UNet~\cite{frunet}, which also ships weights and
is the stronger paper on paper. We moved away from it after reading its code:
it pins an old PyTorch, an old \texttt{timm} that imports a module deleted from
current PyTorch, a Python-2 era helper package, and it calls a YAML function
that no longer exists. That is four fixes before the first epoch, and none of
them teach us anything about vessels. FR-UNet stays as our backup plan.

\begin{table}[!t]
  \caption{Reproducibility check for the selected base paper.}
  \label{tab:checklist}
  \centering
  \renewcommand{\arraystretch}{1.18}
  \setlength{\tabcolsep}{5pt}
  \begin{tabular}{@{} l p{5.4cm} @{}}
    \toprule
    \textbf{Check item} & \textbf{Team's check}                                   \\
    \midrule
    Dataset             & DRIVE, 40 images, already inside the repository         \\
    Official code       & Yes; official PyTorch, MICCAI 2021                      \\
    Protocol            & Fixed 20/20 split supplied by the dataset authors       \\
    Free GPU            & Yes; 20 epochs, $256\times256$ patches, single GPU      \\
    Metrics             & AUC, $F_1$, sensitivity, specificity                    \\
    Weights             & Yes, \texttt{pretrained/drive\_k8.pth}                  \\
    Expected            & Near the published score; exact match not guaranteed    \\
    \bottomrule
  \end{tabular}
\end{table}

\subsection{Hypothesis}

\textbf{Hypothesis:} replacing the pixel-wise loss with an equal mix of the
centre-line clDice loss (Shit \textit{et al.}, CVPR 2021) and cross-entropy
will raise sensitivity by 2 to 5 percentage points and $F_1$ by 0.5 to 1.5
points, by punishing breaks in thin branches far more heavily than a slightly
thick edge, while lowering specificity by no more than 1 point. The network,
the data, the split, the pseudo-labels, the optimiser and the number of epochs
will all stay fixed.

\subsection{Experiment and Ablation Plan}

The main experiment compares the reproduced baseline against the same model
trained with the new loss. The ablation will test: (1)~the published weights,
as a check on our setup; (2)~our retrained baseline with the original loss;
(3)~the same model trained on the original DRIVE labels, without SGL's cleaned
labels; and (4)~the cleaned labels together with the clDice loss. Row~(3)
matters because SGL already repairs missing thin vessels by fixing the labels,
so it is possible that little is left for our loss to win. If that is what we
find, the ablation shows it clearly instead of leaving it as a guess.

We will report $F_1$, sensitivity, specificity, AUC and IoU, and we will treat
any $F_1$ change below 0.01 as noise unless it repeats across seeds, because
the test set is only twenty images. Failure analysis will place the baseline
and the new model side by side on the same images, count broken branches, and
look at where errors gather: faint thin vessels, the bright optic disc, and
crossings where two vessels meet.

% =============================================================
\section{Compact Comparison}

Table~\ref{tab:comparison} places the ten reviewed studies side by side.

\begin{table}[!t]
  \caption{Comparison of the ten reviewed studies on DRIVE. Scores are as
  reported by their own authors under differing protocols, so the table is a
  guide and not a ranking.}
  \label{tab:comparison}
  \centering
  \renewcommand{\arraystretch}{1.18}
  \setlength{\tabcolsep}{4pt}
  \begin{tabular}{@{} l l p{2.5cm} r c @{}}
    \toprule
    \textbf{Study} & \textbf{Venue} & \textbf{Main method} &
    \textbf{$F_1$} & \textbf{Code} \\
    \midrule
    SGL 2021~\cite{sgl}          & MICCAI       & Cleaned labels from a model group & --    & Yes \\
    FR-UNet 2022~\cite{frunet}   & JBHI         & Full-size path, 2-threshold step   & 83$^{*}$ & Yes \\
    DSCNet 2023~\cite{dscnet}    & ICCV         & Bending conv.\ plus shape loss     & 82.06 & Yes \\
    MRC-Net 2023~\cite{mrcnet}   & Neural Netw. & Multi-scale, discriminator         & --    & No  \\
    MDFI-Net 2024~\cite{mdfinet} & arXiv        & Deformable pulse-coupled block     & --    & No  \\
    Neuron prog.\ 2024~\cite{neuronprog} & arXiv & Search over neuron types          & --    & No  \\
    LVS-Net 2024~\cite{lvsnet}   & arXiv        & Light multi-scale encoder          & 86.44 & No  \\
    LFRA-Net 2025~\cite{lfranet} & arXiv        & Focal and region attention         & 84.28 & Yes \\
    MM-UNet 2025~\cite{mmunet}   & IEEE BIBM    & Mamba state-space layers           & $+1.64^{\dagger}$ & Yes \\
    SA-UNetv2 2026~\cite{saunetv2} & IEEE ISBI  & Attention on all skips             & 82.82 & Yes \\
    \bottomrule
  \end{tabular}

  \vspace{3pt}
  {\footnotesize $^{*}$Sources disagree (82.88--83.16); the exact target must be
  read from the paper's own table. $^{\dagger}$Reported as a gain over compared
  methods, not an absolute score. ``--'' means the value was not extracted at
  survey time.}
\end{table}

% =============================================================
\section{Conclusion}

The reviewed papers move from plain U-Net towards full-resolution designs,
cleaner labels, shape-aware convolutions, attention and state-space layers, and
very small models. Scores have largely flattened while model sizes have shrunk,
so a fair and repeatable test now matters more than another point of $F_1$. We
picked SGL as our base paper because everything needed to run it sits in one
repository and because its loss can be swapped from the command line. We will
first reproduce its baseline, then test one paper-backed change aimed at vessel
connectivity. The fixed split, the ablation and the failure analysis together
will show whether that change is worth keeping.

% =============================================================
\section*{Acknowledgment}
The authors thank the open-source communities behind PyTorch and the public
release of the DRIVE benchmark.

% =============================================================
\begin{thebibliography}{00}

\bibitem{frunet}
W.~Liu \textit{et al.},
``Full-resolution network and dual-threshold iteration for retinal vessel and
coronary angiograph segmentation,''
\textit{IEEE J.\ Biomed.\ Health Inform.}, vol.~26, no.~9, pp.~4623--4634,
2022, doi: 10.1109/JBHI.2022.3188710.

\bibitem{mdfinet}
Y.~Dong and X.~Deng,
``MDFI-Net: Multiscale differential feature interaction network for accurate
retinal vessel segmentation,''
\textit{arXiv preprint arXiv:2410.15444}, 2024,
doi: 10.48550/arXiv.2410.15444.

\bibitem{sgl}
Y.~Zhou, H.~Yu, and H.~Shi,
``Study group learning: Improving retinal vessel segmentation trained with
noisy labels,''
in \textit{Proc.\ MICCAI}, LNCS vol.~12901, pp.~57--67, 2021,
doi: 10.1007/978-3-030-87193-2\_6.

\bibitem{mrcnet}
T.~M.~Khan, S.~S.~Naqvi, A.~Robles-Kelly, and I.~Razzak,
``Retinal vessel segmentation via a multi-resolution contextual network and
adversarial learning,''
\textit{Neural Networks}, vol.~165, pp.~310--320, 2023,
doi: 10.1016/j.neunet.2023.05.029.

\bibitem{dscnet}
Y.~Qi, Y.~He, X.~Qi, Y.~Zhang, and G.~Yang,
``Dynamic snake convolution based on topological geometric constraints for
tubular structure segmentation,''
in \textit{Proc.\ IEEE/CVF ICCV}, pp.~6047--6056, 2023,
doi: 10.1109/ICCV51070.2023.00558.

\bibitem{saunetv2}
C.~Guo, A.~N.~Christensen, A.~B.~Dahl, Y.~Yi, and M.~R.~Hannemose,
``SA-UNetv2: Rethinking spatial attention U-Net for retinal vessel
segmentation,''
in \textit{Proc.\ IEEE ISBI}, pp.~1--5, 2026,
doi: 10.1109/ISBI61048.2026.11515447.

\bibitem{mmunet}
J.~Liu \textit{et al.},
``MM-UNet: Morph Mamba U-shaped convolutional networks for retinal vessel
segmentation,''
in \textit{Proc.\ IEEE BIBM}, 2025, doi: 10.1109/BIBM66473.2025.11357047.

\bibitem{lvsnet}
M.~Mehmood, S.~Iqbal, T.~M.~Khan, I.~Spence, and M.~Fahim,
``LVS-Net: A lightweight vessels segmentation network for retinal image
analysis,''
\textit{arXiv preprint arXiv:2412.05968}, 2024,
doi: 10.48550/arXiv.2412.05968.

\bibitem{lfranet}
M.~Mehmood, S.~Iqbal, T.~M.~Khan, I.~Spence, and M.~Fahim,
``LFRA-Net: A lightweight focal and region-aware attention network for retinal
vessel segmentation,''
\textit{arXiv preprint arXiv:2509.11811}, 2025,
doi: 10.48550/arXiv.2509.11811.

\bibitem{neuronprog}
T.~Wu, R.~Min, P.~Song, H.~Guo, T.~Zeng, and F.-L.~Fan,
``Retinal vessel segmentation via neuron programming,''
\textit{arXiv preprint arXiv:2411.11110}, 2024,
doi: 10.48550/arXiv.2411.11110.

\end{thebibliography}

\end{document}
